\documentclass[11pt]{article}
\usepackage[T1]{fontenc}
\usepackage{lmodern}
\usepackage{amsmath,amssymb,amsthm}
\usepackage[margin=1in]{geometry}
\usepackage[hidelinks]{hyperref}
\newtheorem{proposition}{Proposition}
\newtheorem{corollary}{Corollary}
\title{A Flat Maximum in the Quantum Factorial:\\Disproof of Conjecture 00000007683}
\author{\texttt{ziangni-sys}\\\small Prepared with OpenAI Codex}
\date{4 October 2026}
\begin{document}
\maketitle

\begin{abstract}
The coefficients of $[3]_q!$ are $(1,2,2,1)$. In particular, the two adjacent
middle coefficients form a flat maximum. This contradicts the conjecture's
assertion that the coefficient sequence is flat-free for every fixed $n$.
The defining polynomial multiplication, all coefficients, and the resulting
contradiction are verified by a self-contained Lean 4 project.
\end{abstract}

\section{Statement and scope}
The source defines
\[
 [n]_q! = \prod_{i=1}^{n}(1+q+\cdots+q^{i-1})
         = \sum_{m\geq 0} c_{n,m}q^m
\]
and asserts, in particular, that for each fixed $n$ the coefficient sequence
is ``flat-free and log-concave.'' The Chinese statement likewise asserts the
absence of a flat portion. The additional assertions concern the location
and asymptotic size of a peak.

We use ``flat-free'' to mean the absence of equal adjacent entries in the
finite coefficient sequence. Our example has the stronger property that
two adjacent entries are both global maxima. Thus it also contradicts the
absence of a flat peak. No trailing zero coefficients, limiting regime, or
choice of numerical approximation is involved.

\section{Exact counterexample}
\begin{proposition}
For $n=3$,
\[
 [3]_q! = 1+2q+2q^2+q^3.
\]
The maximum coefficient is $2$, attained at both degrees $1$ and $2$.
\end{proposition}
\begin{proof}
Directly from the stated definition,
\begin{align*}
 [3]_q! &= (1)(1+q)(1+q+q^2)\\
 &= (1+q+q^2)+(q+q^2+q^3)\\
 &= 1+2q+2q^2+q^3.
\end{align*}
Consequently $c_{3,0}=c_{3,3}=1$, $c_{3,1}=c_{3,2}=2$, and all
coefficients at degrees at least $4$ are zero. In particular, the equal
coefficients at degrees $1$ and $2$ are positive, inside the support,
and greater than their two outside neighbors.
\end{proof}

\begin{corollary}
Conjecture 00000007683, as stated, is false.
\end{corollary}
\begin{proof}
Its universal flat-free assertion includes $n=3$, whereas the proposition
exhibits a flat peak at that value. A false necessary clause makes the
conjunction of all asserted clauses false.
\end{proof}

For clarity, the example is log-concave: both interior inequalities are
$2^2\geq 1\cdot 2$. It is the flat-free assertion that fails. Neither a
disproof of log-concavity nor an analysis of the peak asymptotic is claimed
or required here.

\section{Lean certificate and semantic correspondence}
The accompanying project uses Lean 4.19.0 and only its bundled
\texttt{Std} library. A polynomial is represented by its list of natural
number coefficients in increasing degree order. These coefficients are
nonnegative integers, as in the defining product above.

\begin{itemize}
\item \texttt{add} is coefficientwise addition.
\texttt{mul} implements the convolution identity
$(a+XP)Q=aQ+X(PQ)$.
\item \texttt{eval\_add}, \texttt{eval\_scale}, and
\texttt{eval\_mul} establish the corresponding polynomial evaluation
identities for every natural-number argument, not merely the example.
\item \texttt{qInteger n} is the list of $n$ ones, representing
$1+q+\cdots+q^{n-1}$.
\texttt{qFactorial} recursively multiplies these lists, beginning with
the empty product $1$.
\item \texttt{qFactorial\_three} computes the entire coefficient list
by definitional equality. The lemma \path{coefficient_formula} identifies
the coefficient at every natural-number degree, including those beyond
the support.
\item \texttt{flat\_maximum} proves the existence of an adjacent
equal pair inside the finite coefficient sequence and that both entries
are global maxima.
\item \texttt{conjecture\_flatFree\_false} refutes the universal
flat-free clause even when restricted to $n\geq 3$.
\texttt{conjecture\_false} establishes that conjoining arbitrary remaining
clauses with this necessary clause is impossible.
\end{itemize}

The final contradiction theorems have no axiom dependencies. Some auxiliary
evaluation and maximum lemmas use standard Lean logical axioms; their exact
dependencies are printed by the project. There are no admitted proofs,
external proof oracles, or added axioms. From the \texttt{lean} directory,
run \texttt{lake build} and
\texttt{lake env lean -DwarningAsError=true Main.lean}.

The optional Python check independently multiplies the coefficient lists
and enumerates the inversion numbers of all six permutations of three
elements. It is supplementary evidence, not a premise of the Lean proof.

\section{Source and attribution}
The statement is \texttt{conjectures/00000007683.md} in
\href{https://github.com/The-Last-Math-Competition/The-Last-Math-Competition}
{The Last Math Competition}, at commit
\begin{center}
\small\texttt{4509328ceb3b63c07d05bc4c42c4f12e8fa3c94f}.
\end{center}
The complete English and Chinese source text is reproduced in
\texttt{SOURCE.md}. This submission was prepared by OpenAI Codex for
\texttt{ziangni-sys} and is offered for independent review.
\end{document}
